This thesis addressed the problem of automated FAIRness assessment
for research software repositories, motivated by the observation
that existing tools primarily evaluate structured metadata and
repository-level artefacts but do not directly assess qualitative
characteristics of unstructured software documentation. To address
this gap, a hybrid assessment framework was designed, implemented,
and evaluated, combining a deterministic rule-based component
covering six FAIR4RS-derived structural indicators with an LLM-based
component evaluating six documentation-quality dimensions of
repository README content.

The framework was applied to a dataset of 188 research software
repositories collected from GitHub across ten research domains, with
a manually annotated ground truth corpus of 30 repositories
constructed to evaluate both assessment components. The rule-based
component was shown to reproduce its operational definitions
faithfully against the finalised ground truth and revealed substantial
variability in the structural FAIR-related indicators assessed across
the full dataset.

The LLM-based component was evaluated using two model configurations:
a locally hosted Llama~3.1 8B model and
\texttt{Qwen3.6-35B-A3B-MLX-8bit}, accessed through a
university-hosted API. The two configurations produced substantially
different levels of agreement with the primary human annotator.
Llama achieved a macro-average weighted Cohen's Kappa of 0.329,
whereas Qwen achieved 0.612. Against the second annotator, the
corresponding values were 0.275 and 0.557. These results show that,
within the evaluated configurations, the reliability of qualitative
LLM-based documentation assessment varies with both the model
configuration and the documentation dimension being assessed.

In relation to the two research questions guiding this work, the
results indicate that LLM-based assessment can complement rule-based
FAIRness assessment by evaluating qualitative documentation
characteristics that are not represented by the structural indicators
considered in this work. The associations between the structural and
documentation totals were weak and non-significant for both Llama
($\rho=0.175$, $p=0.356$) and Qwen
($\rho=0.134$, $p=0.480$). These findings are consistent with the two
components contributing different information rather than providing
redundant measurements of the same repository characteristics,
although they do not establish statistical independence. Their
combination therefore broadens the information represented in an
automated research-software FAIRness assessment while preserving the
distinction between structural and qualitative evidence (RQ1).

For RQ2, the two-model evaluation shows that LLM-based assessment
reliability is neither uniform across documentation dimensions nor
stable across model configurations. Qwen showed strong agreement for
installation, purpose, and contribution information, while dependency
documentation and examples remained more difficult. Examples (D5)
showed systematic underscoring under both model configurations, while
Qwen additionally showed pronounced underscoring for dependency
documentation (D4). The contrast for purpose description (D3) was
particularly notable: Llama achieved $\kappa=-0.071$ against the
primary annotator, whereas Qwen achieved $\kappa=0.789$. This
demonstrates that a weak result for an individual dimension under one
model configuration should not automatically be interpreted as an
inherent limitation of LLM-based assessment for that dimension.

The reliability analyses also identified limitations that persisted
despite the stronger Qwen agreement. Model-reported confidence did
not provide a dependable reliability signal. Llama's confidence
levels were not monotonically associated with observed correctness,
while Qwen assigned HIGH confidence to all 180 dimension-level
assessments, preventing confidence from discriminating between more
and less reliable predictions. At the repository level, no
statistically significant relationship was observed between README
length and model--human scoring disagreement for either model
configuration. These findings reinforce the need to evaluate
LLM-generated assessments empirically and to retain transparent
criterion-level evidence rather than relying on model confidence
alone.

The primary contribution of this thesis is a working hybrid FAIRness
assessment framework for research software repositories, together
with a web-based interface for inspecting dataset-level results,
repository-level profiles, evaluation results, model comparisons,
and live assessments of GitHub repositories. The implemented
prototype supports both evaluated model configurations for live
documentation assessment: the local Llama configuration through
Ollama and the Qwen configuration through the university-hosted
KISTE service. The framework preserves structural and
documentation-based evidence as separate, interpretable dimensions
rather than collapsing them into a single overall FAIRness score.

Its empirical evaluation against manual ground truth,
human inter-annotator agreement, and existing FAIR assessment
approaches characterises both the additional assessment coverage and
the reliability limitations of the proposed approach. The second
model evaluation further demonstrates that the choice between the
evaluated model configurations materially affected the observed
reliability of qualitative documentation assessment.
At the same time, recurring difficulties for particular dimensions,
especially examples and dependency documentation, show that stronger
aggregate model agreement does not eliminate criterion-specific
failure patterns.

Taken together, this work demonstrates that combining rule-based and
LLM-based assessment can broaden the scope of automated FAIRness
evaluation for research software by bringing structural and
documentation-based evidence into a common assessment framework.
However, the results also show that LLM-based assessment should not
be treated as uniformly reliable: its performance depends on the
model configuration, the documentation dimension, and the available
evidence. The hybrid approach therefore provides the greatest value
when LLM assessment is used as a complement to deterministic checks,
with explicit criterion-level results that remain open to human
interpretation. In this way, the framework provides a practical and
empirically evaluated basis for extending automated FAIRness
assessment beyond structural repository checks while preserving the
transparency needed for responsible interpretation of LLM-generated
assessments.